% ============================================================
% Preamble untuk Buku Ajar OBE - Pemrograman Game
% Menggabungkan OBE environments + C# listings
% ============================================================

% --- Encoding dan Bahasa ---
\usepackage[utf8]{inputenc}
\usepackage[T1]{fontenc}
\usepackage[indonesian]{babel}
\usepackage{lmodern}
\usepackage{times}

% --- Layout Halaman ---
\usepackage{geometry}
\geometry{a4paper, left=3cm, right=2cm, top=2.5cm, bottom=2cm, headheight=15pt}

% --- Grafik dan Tabel ---
\usepackage{graphicx}
\graphicspath{{figures/}{../figures/}}
\usepackage{float}
\usepackage{booktabs}
\usepackage{longtable}
\usepackage{array}
\usepackage{multirow}
\usepackage{xcolor}
\usepackage{amssymb}
\usepackage{amsmath,amsthm}

% --- TikZ untuk Diagram ---
\usepackage{tikz}
\usetikzlibrary{shapes,arrows,positioning,calc}
\tikzset{
  block/.style={rectangle, draw, fill=blue!20, minimum height=1.2em, minimum width=2.5em, align=center, rounded corners=2pt},
  line/.style={draw},
  class/.style={rectangle, draw=black, fill=blue!10, minimum width=3cm, minimum height=1cm, text centered, font=\bfseries},
  object/.style={rectangle, rounded corners, draw=black, fill=green!10, minimum width=3cm, minimum height=1cm, text centered},
  arrow/.style={thick, ->, >=stealth}
}

% --- List dan Enumerasi ---
\usepackage{enumitem}
\setlist[itemize]{leftmargin=*,topsep=0.5em}
\setlist[enumerate]{leftmargin=*,topsep=0.5em}

% --- Subfiles untuk Modular Structure ---
\usepackage{subfiles}
\providecommand{\ifSubfilesClassLoaded}[2]{\IfSubfilesClassLoadedTF{#1}{#2}}

% --- Code Listings untuk C# ---
\usepackage{listings}
\definecolor{codegray}{rgb}{0.5,0.5,0.5}
\definecolor{codebg}{rgb}{0.95,0.95,0.95}
\lstdefinestyle{csharpstyle}{
  language=[Sharp]C,
  basicstyle=\ttfamily\small,
  keywordstyle=\color{blue}\bfseries,
  commentstyle=\color{green!60!black},
  stringstyle=\color{red},
  numbers=left,
  numberstyle=\tiny\color{codegray},
  stepnumber=1,
  numbersep=10pt,
  backgroundcolor=\color{codebg},
  showspaces=false,
  showstringspaces=false,
  frame=single,
  tabsize=2,
  captionpos=b,
  breaklines=true,
  breakatwhitespace=false,
  escapeinside={(*@}{@*)},
  xleftmargin=15pt,
  xrightmargin=5pt
}
\renewcommand{\lstlistlistingname}{Daftar Kode Program}
\renewcommand{\lstlistingname}{Kode Program}
\lstset{style=csharpstyle}

% --- Theorems untuk Definisi, Contoh, Catatan ---
\newtheorem{definisi}{Definisi}[section]
\newtheorem{contohtheorem}{Contoh}[section]
\newtheorem{tips}{Tips}[section]
\newenvironment{learningoutcomes}{%
  \par\noindent\textbf{Learning Outcomes:}\par\vspace{0.3em}%
  \begin{itemize}[leftmargin=*, itemsep=0.2em]%
}{%
  \end{itemize}%
}

% --- Boxes untuk Highlight ---
\usepackage{tcolorbox}
\tcbuselibrary{skins,breakable}
\newtcolorbox{catatan}{
  colback=yellow!5!white,
  colframe=yellow!75!black,
  fonttitle=\bfseries,
  title=Catatan,
  breakable
}
\newtcolorbox{contoh}{
  colback=green!5!white,
  colframe=green!75!black,
  fonttitle=\bfseries,
  title=Contoh,
  breakable
}
\newtcolorbox{konsep}{
  colback=blue!5!white,
  colframe=blue!75!black,
  fonttitle=\bfseries,
  title=Konsep Penting,
  breakable
}

% --- Custom Environments untuk OBE ---
\newenvironment{subcpmk}{
  \vspace{0.5cm}
  \noindent\colorbox{blue!10}{\parbox{\dimexpr\textwidth-2\fboxsep}{
    \textbf{\large Sub-CPMK yang Dicakup dalam Bab Ini:}
  }}
  \vspace{0.3cm}
  \begin{itemize}[leftmargin=*, itemsep=5pt]
}{
  \end{itemize}
  \vspace{0.5cm}
}

\newenvironment{aktivitas}{
  \vspace{0.5cm}
  \noindent\colorbox{green!10}{\parbox{\dimexpr\textwidth-2\fboxsep}{
    \textbf{\large Aktivitas Pembelajaran}
  }}
  \vspace{0.3cm}
  \begin{enumerate}[leftmargin=*, itemsep=8pt]
}{
  \end{enumerate}
  \vspace{0.5cm}
}

\newenvironment{latihan}{
  \vspace{0.5cm}
  \noindent\colorbox{orange!10}{\parbox{\dimexpr\textwidth-2\fboxsep}{
    \textbf{\large Latihan dan Refleksi}
  }}
  \vspace{0.3cm}
  \begin{enumerate}[leftmargin=*, itemsep=8pt]
}{
  \end{enumerate}
  \vspace{0.5cm}
}

\newenvironment{asesmen}{
  \vspace{0.5cm}
  \noindent\colorbox{red!10}{\parbox{\dimexpr\textwidth-2\fboxsep}{
    \textbf{\large Asesmen (Evaluasi Kinerja)}
  }}
  \vspace{0.3cm}
}{
  \vspace{0.5cm}
}

\newenvironment{checklist}{
  \vspace{0.5cm}
  \noindent\colorbox{purple!10}{\parbox{\dimexpr\textwidth-2\fboxsep}{
    \textbf{\large Checklist Pencapaian Kompetensi}
  }}
  \vspace{0.3cm}
  \noindent\textit{Centang item berikut setelah Anda yakin telah menguasainya:}
  \vspace{0.2cm}
  \begin{itemize}[leftmargin=*, itemsep=5pt, label=$\square$]
}{
  \end{itemize}
  \vspace{0.5cm}
}

\newenvironment{rangkuman}{
  \vspace{0.5cm}
  \noindent\colorbox{gray!10}{\parbox{\dimexpr\textwidth-2\fboxsep}{
    \textbf{\large Rangkuman}
  }}
  \vspace{0.3cm}
}{
  \vspace{0.5cm}
}

% --- Custom Commands ---
\newcommand{\code}[1]{\texttt{#1}}
\newcommand{\class}[1]{\texttt{\textbf{#1}}}
\newcommand{\method}[1]{\texttt{#1()}}
\newcommand{\keyword}[1]{\texttt{\textbf{#1}}}

% --- Hyperlinks ---
\usepackage{xurl}
\usepackage{hyperref}
\hypersetup{
  colorlinks=true,
  linkcolor=blue,
  urlcolor=blue,
  citecolor=blue,
  bookmarksnumbered=true,
  pdfstartview=FitH,
  pdftitle={Buku Ajar Pemrograman Game},
  pdfsubject={Game Development dengan Unity}
}

% --- Header dan Footer ---
\usepackage{fancyhdr}
\pagestyle{fancy}
\fancyhf{}
\fancyhead[LE,RO]{\thepage}
\fancyhead[RE]{\leftmark}
\fancyhead[LO]{\rightmark}
\renewcommand{\headrulewidth}{0.4pt}

% --- Spacing ---
\setlength{\parskip}{0.5em}
\setlength{\parindent}{1.5em}
\renewcommand{\baselinestretch}{1.5}
\setcounter{tocdepth}{2}
\setcounter{secnumdepth}{3}
